\begin{pastexam}[label={px:ccu-109-csb-02}]{中正資工所乙組 109 學年度 第 5 題}
（是非）

(20\%) Carefully examine the following twenty sentences and give your answer, true (T) or false
(F), for each sentence. Two points are given for each correct answer, whereas \textbf{three points
are deducted} for each incorrect answer. Therefore, please do not guess at the answer if you are
unsure since \textbf{the expected value is negative.} Note that the deduction for incorrect
responses only affects the 20 points of this question.
\begin{enumerate}[label=(\alph*)]
\item The fields of a user-defined data type ``union'' in C share their memory space.
\item The maximum number of nodes in a binary tree of depth 10 is 1024.
\item We can invert a linked list in place (i.e., no need to use additional memory with a size
proportional to the elements in the linked list).
\item A heap (i.e., priority queue) is better than a binary search tree when we examine whether
an element is stored in the data set.
\item Dijkstra algorithm cannot deal with the graph with negative edge weights.
\item In most cases, prime numbers are good divisors for implementing a division-based hash
function.
\item A balanced binary search tree is always the best binary search tree for storing the records
even if the identifiers are with different searching probabilities.
\item Quick sort is faster than insertion sort even if the input size is small.
\item A stack is a first-in-first-out data structure.
\item A user-defined ``struct'' in C creates a data type that can be used to group items of
possibly different types into a single type.
\end{enumerate}
\end{pastexam}

\begin{pastsoln}
答案：
\begin{center}
\begin{tabular}{*{10}{c}}
\toprule
(a) & (b) & (c) & (d) & (e) & (f) & (g) & (h) & (i) & (j) \\
\midrule
T & F & T & F & T & T & F & F & F & T \\
\bottomrule
\end{tabular}
\end{center}
題目說「twenty sentences」，實際只有 (a)--(j) 十句，每句 $2$ 分，合起來正好 $20$ 分，「twenty」
應該是筆誤。 (a)、(j) 考 C 語言，不在資結書的範圍，這裡簡單說明。

\textbf{(a) T}。 \texttt{union} 的所有欄位從同一個位址開始存放，共用同一塊記憶體，大小是最大的那個
欄位（再加上對齊）。 同一時間只有最後寫入的那個欄位有意義。

\textbf{(b) F}。 \textbf{depth} 有兩種常見的定義，兩種都得不到 $1024$：
\begin{itemize}
\item depth 指層數、根在第 $1$ 層（Horowitz 課本的定義）：最多 $2^{10} - 1 = 1023$ 個節點。
\item depth 指根到最深節點的邊數（資結書第~7.1~節的深度）：共 $11$ 層，最多
$2^{11} - 1 = 2047$ 個。
\end{itemize}
第 $i$ 層最多 $2^{i-1}$ 個節點，總數是 $1 + 2 + \cdots$ 這種等比級數的和，一定是
$2^k - 1$ 的形式（資結書定理~7.1），不可能是 $1024 = 2^{10}$。 $1024$ 只是「深度為 $10$
的那一層」最多的節點數，不是整棵樹的。

\textbf{(c) T}。 就地反轉只要三個指標 $\mathit{prev}$、$\mathit{cur}$、$\mathit{next}$：從頭走到尾，
把每個節點的 $\mathit{next}$ 改指向前一個節點（資結書虛擬碼~5.12）。 額外空間是 $O(1)$，
和串列長度無關。

\textbf{(d) F}。 堆積只保證父節點不小於（或不大於）子節點，兄弟之間、左右子樹之間沒有大小關係，
找一個任意的鍵沒辦法決定往哪邊走，最壞要看過全部 $n$ 個元素，$O(n)$。 二元搜尋樹每次比較就能
決定往左或往右，花 $O(h)$；平衡時 $h = O(\lg n)$。 判斷「某個元素在不在」時二元搜尋樹比較好；
堆積擅長的是取最大（最小）值（資結書第~8.3~節）。

\textbf{(e) T}。 Dijkstra 演算法每次取出距離最小的頂點後就把它的距離定下來，不再更新；這要靠
「之後的路徑只會越走越長」，邊權重是負的時候不成立（資結書第~25.2.5~節）。 例如三條邊
$1 \to 2$（權重 $3$）、$1 \to 3$（$5$）、$3 \to 2$（$-4$），從 $1$ 出發：Dijkstra 先取出 $2$，
定下 $d[2] = 3$；但 $1 \to 3 \to 2$ 的長度是 $1$。 所以有負邊的圖不能用 Dijkstra，要改用
Bellman--Ford。 有負邊時 Dijkstra 偶爾也會碰巧算對，但不保證，一般就說它不能處理負邊。

\textbf{(f) T}。 除法雜湊 $h(k) = k \bmod m$ 的 $m$ 通常選質數，而且不要太接近 $2$ 的次方
（資結書第~12.1~節）。 $m$ 和鍵常見的間隔（例如全是偶數、全是 $10$ 的倍數）有公因數時，鍵只會
落在一部分的格子；$m$ 是質數就不會有這個問題。 題目說「in most cases」，對。

\textbf{(g) F}。 各個鍵被搜尋的機率不同時，最好的二元搜尋樹要把常被搜尋的鍵放在靠近根的地方，
不一定平衡，這就是最佳二元搜尋樹（資結書第~18.4~節）。 例如鍵 $1, 2, 3$ 的搜尋機率是
$0.8, 0.1, 0.1$：平衡的樹以 $2$ 為根，期望比較次數 $0.8 \cdot 2 + 0.1 \cdot 1 + 0.1 \cdot 2 = 1.9$；
以 $1$ 為根、$2$ 接右邊、$3$ 再接右邊的斜樹只要 $0.8 \cdot 1 + 0.1 \cdot 2 + 0.1 \cdot 3 = 1.3$。

\textbf{(h) F}。 $n$ 很小時，插入排序的常數小、沒有遞迴呼叫的負擔，比快速排序快（資結書
第~14.2~節）。 實務上的快速排序在子陣列小到某個長度以下時改用插入排序，就是這個原因。

\textbf{(i) F}。 堆疊是後進先出（last-in-first-out, LIFO）；先進先出（FIFO）的是佇列（資結書
第~6.1、6.2~節）。

\textbf{(j) T}。 \texttt{struct} 把幾個可能不同型別的欄位組成一個新的型別，每個欄位各有自己的
記憶體位置。 這正是它和 (a) 的 \texttt{union} 不同的地方。
\end{pastsoln}
